\needspace{20\baselineskip}
\subsection*{Dot products and norms}

\begin{problem}\label{ex:01:orthogonality-transitivity}
Give vectors $x,y,z$ for which $x\perp y$ and $y\perp z$, but
$x$ is not perpendicular to $z$.
\end{problem}
\begin{solution}
Take $x=z=(1,0)$ and $y=(0,1)$. Then
$x\cdot y=y\cdot z=0$, whereas $x\cdot z=1$.
\end{solution}

\begin{problem}\label{ex:01:parallelogram-polarization}
For $x,y\in\R^n$, prove
\[
\norm{x+y}^2+\norm{x-y}^2=2\norm{x}^2+2\norm{y}^2,
\qquad
\norm{x+y}^2-\norm{x-y}^2=4x\cdot y.
\]
Interpret the first identity as a parallelogram law.
\end{problem}
\begin{solution}
The dot-product laws give
\[
\norm{x+y}^2=\norm{x}^2+2x\cdot y+\norm{y}^2,
\qquad
\norm{x-y}^2=\norm{x}^2-2x\cdot y+\norm{y}^2.
\]
Add and subtract these equations. The parallelogram with vertices
$0,x,x+y,y$ has diagonals of lengths $\norm{x+y}$ and $\norm{x-y}$.
The first identity says that the sum of their squared lengths equals
the sum of the squared lengths of the four sides.
\end{solution}

\begin{problem}\label{ex:01:dot-cancellation}
Give three nonzero vectors $a,b,c$ such that $a\cdot b=a\cdot c$
but $b\ne c$.
\end{problem}
\begin{solution}
Take $a=(1,0)$, $b=(1,1)$, and $c=(1,2)$. Both dot products are $1$,
although $b\ne c$.
\end{solution}

\begin{problem}\label{ex:01:perpendicular-sum}
Let $u_1,\ldots,u_r\in\R^n$ be nonzero and pairwise perpendicular.
If $c_1u_1+\cdots+c_ru_r=0$, prove that $c_1=\cdots=c_r=0$.
\end{problem}
\begin{solution}
For each $j$, take the dot product with $u_j$:
\[
0=\sum_{i=1}^r c_i(u_i\cdot u_j)=c_j\norm{u_j}^2.
\]
Since $u_j\ne0$, its squared norm is positive, so $c_j=0$.
\end{solution}

\begin{problem}\label{ex:01:numerical-projection}
For $a=(2,-1,5)$ and $b=(-1,1,1)$, find the projection of $a$ along
$b$ and the projection of $b$ along $a$.
\end{problem}
\begin{solution}
Here $a\cdot b=-2-1+5=2$, $\norm{a}^2=30$, and
$\norm{b}^2=3$. Thus the projections are, respectively,
\[
\frac{a\cdot b}{\norm{b}^2}b
=\left(-\frac23,\frac23,\frac23\right),
\qquad
\frac{b\cdot a}{\norm{a}^2}a
=\left(\frac{2}{15},-\frac{1}{15},\frac13\right).
\]
\end{solution}

\begin{problem}\label{ex:01:triangle-equality}
Use the Schwarz inequality to prove that, for $x,y\in\R^n$,
\[
\norm{x+y}\leqslant\norm{x}+\norm{y}.
\]
Prove that equality holds exactly when $y=0$ or $x=ky$ for some $k\geqslant0$.
\end{problem}
\begin{solution}
By Schwarz,
\[
\norm{x+y}^2
=\norm{x}^2+2x\cdot y+\norm{y}^2
\leqslant\norm{x}^2+2\norm{x}\norm{y}+\norm{y}^2
=(\norm{x}+\norm{y})^2.
\]
Taking nonnegative square roots proves the inequality. Equality requires
$x\cdot y=\norm{x}\norm{y}$. If $y=0$, this holds. If $y\ne0$,
the equality condition in Schwarz requires $x=ky$ for a real $k$.
Then $x\cdot y=k\norm{y}^2$ and
$\norm{x}\norm{y}=|k|\norm{y}^2$, which are equal exactly when $k\geqslant0$.
\end{solution}

\needspace{20\baselineskip}
\subsection*{Lines and planes}

\begin{problem}\label{ex:01:segment-proportion}
Let $x\ne y$ and $0\leqslant\lambda\leqslant1$. Show that
$z=(1-\lambda)x+\lambda y$ divides the segment from $x$ to $y$
in the proportion $\lambda:(1-\lambda)$.
\end{problem}
\begin{solution}
The representation $z=x+\lambda(y-x)$ places $z$ on the segment.
Moreover,
\[
\norm{z-x}=\norm{\lambda(y-x)}=\lambda\norm{y-x},
\qquad
\norm{y-z}=(1-\lambda)\norm{y-x}.
\]
Since $\norm{y-x}>0$, these distances have the stated proportion,
including $z=x$ at $\lambda=0$ and $z=y$ at $\lambda=1$.
\end{solution}

\begin{problem}\label{ex:01:plane-normal}
Find the equation of the plane perpendicular to $(1,-1,3)$ and passing
through $(4,2,-1)$.
\end{problem}
\begin{solution}
A point $(x_1,x_2,x_3)$ lies on the plane exactly when
\[
(1,-1,3)\cdot(x_1-4,x_2-2,x_3+1)=0.
\]
Thus its equation is $x_1-x_2+3x_3=-1$.
\end{solution}

\begin{problem}\label{ex:01:perpendicular-slopes}
Write down vectors perpendicular to the lines $y=mx+b$ and $y=m'x+c$.
Show that these vectors are perpendicular to each other exactly when
$mm'=-1$.
\end{problem}
\begin{solution}
The equations $-mx+y=b$ and $-m'x+y=c$ give nonzero normal vectors
$(-m,1)$ and $(-m',1)$. Their dot product is $mm'+1$,
which vanishes exactly when $mm'=-1$.
\end{solution}

\begin{problem}\label{ex:01:line-plane-intersection}
\begin{enumerate}
\item Find the intersection of the line through $(1,3,5)$ with direction
$(-2,1,1)$ and the plane $2x+3y-z=1$.
\item Find the intersection of the plane $3x-4y+z=2$ and the line through
$(1,2,-1)$ perpendicular to that plane.
\end{enumerate}
\end{problem}
\begin{solution}
(a) Substitute $(x,y,z)=(1-2t,3+t,5+t)$ into the plane equation:
\[
2(1-2t)+3(3+t)-(5+t)=6-2t=1.
\]
Hence $t=5/2$ and the intersection is $(-4,11/2,15/2)$.

(b) The normal vector $(3,-4,1)$ gives the line
$(x,y,z)=(1+3t,2-4t,-1+t)$. Substitution gives
\[
3(1+3t)-4(2-4t)+(-1+t)=-6+26t=2.
\]
Thus $t=4/13$ and the intersection is $(25/13,10/13,-9/13)$.
\end{solution}

\needspace{20\baselineskip}
\subsection*{Matrix multiplication and transpose}

\begin{problem}\label{ex:01:rectangular-associativity}
Let
\[
A=\begin{pmatrix}2&1&-1\\3&1&2\end{pmatrix},\qquad
B=\begin{pmatrix}1&1\\2&0\\3&-1\end{pmatrix},\qquad
c=\begin{pmatrix}1\\3\end{pmatrix}.
\]
Compute $(AB)c$ and $A(Bc)$, showing the intermediate products.
\end{problem}
\begin{solution}
The intermediate products have sizes $2\times2$ and $3\times1$:
\[
AB=\begin{pmatrix}2+2-3&2+1\\3+2+6&3-2\end{pmatrix}
=\begin{pmatrix}1&3\\11&1\end{pmatrix},
\qquad Bc=\begin{pmatrix}4\\2\\0\end{pmatrix}.
\]
Consequently,
\[
(AB)c=\begin{pmatrix}1+9\\11+3\end{pmatrix}
=\begin{pmatrix}10\\14\end{pmatrix},
\qquad
A(Bc)=\begin{pmatrix}8+2\\12+2\end{pmatrix}
=\begin{pmatrix}10\\14\end{pmatrix}.
\]
\end{solution}

\begin{problem}\label{ex:01:reshuffle-order}
On an ordered triple $(a_1,a_2,a_3)$, let $S$ reverse the order and let
$T$ move the last entry to the first position. Show that performing $S$
then $T$ need not give the same result as performing $T$ then $S$.
Do this first with triples, then with matrices acting on column vectors.
\end{problem}
\begin{solution}
Applying $S$ then $T$ gives
$(a_1,a_2,a_3)\mapsto(a_3,a_2,a_1)\mapsto(a_1,a_3,a_2)$.
Applying $T$ then $S$ gives
$(a_1,a_2,a_3)\mapsto(a_3,a_1,a_2)\mapsto(a_2,a_1,a_3)$.
For instance, the triple $(1,2,3)$ gives different results.
The corresponding matrices are
\[
S=\begin{pmatrix}0&0&1\\0&1&0\\1&0&0\end{pmatrix},\qquad
T=\begin{pmatrix}0&0&1\\1&0&0\\0&1&0\end{pmatrix}.
\]
The two orders give
\[
TS=\begin{pmatrix}1&0&0\\0&0&1\\0&1&0\end{pmatrix},\qquad
ST=\begin{pmatrix}0&1&0\\1&0&0\\0&0&1\end{pmatrix},
\]
which are unequal.
\end{solution}

\begin{problem}\label{ex:01:matrix-constructions}
Construct real $2\times2$ matrices satisfying each condition:
\begin{enumerate}
\item $A^2=-I$;
\item $B^2=0$ but $B\ne0$;
\item $CD=-DC$ and $CD\ne0$;
\item $EF=0$, with every entry of both $E$ and $F$ nonzero.
\end{enumerate}
\end{problem}
\begin{solution}
(a) Take $A=\begin{pmatrix}0&-1\\1&0\end{pmatrix}$. Then
$A^2=\begin{pmatrix}-1&0\\0&-1\end{pmatrix}=-I$.

(b) Take $B=\begin{pmatrix}0&1\\0&0\end{pmatrix}$. It is nonzero,
and direct multiplication gives $B^2=\begin{pmatrix}0&0\\0&0\end{pmatrix}$.

(c) Take
\[
C=\begin{pmatrix}1&0\\0&-1\end{pmatrix},\qquad
D=\begin{pmatrix}0&1\\1&0\end{pmatrix}.
\]
Then $CD=\begin{pmatrix}0&1\\-1&0\end{pmatrix}$ and
$DC=\begin{pmatrix}0&-1\\1&0\end{pmatrix}=-CD$.

(d) Take $E=\begin{pmatrix}1&1\\1&1\end{pmatrix}$ and
$F=\begin{pmatrix}1&-1\\-1&1\end{pmatrix}$. Every entry is nonzero,
but each entry of $EF$ is either $1-1$ or $-1+1$, so $EF=0$.
\end{solution}

\begin{problem}\label{ex:01:difference-of-squares}
Two square matrices $A,B$ of the same size \emph{commute} if $AB=BA$.
Prove that
\[
(A+B)(A-B)=A^2-B^2
\]
holds exactly when $A$ and $B$ commute.
\end{problem}
\begin{solution}
Distributivity gives
\[
(A+B)(A-B)=A^2-AB+BA-B^2.
\]
This equals $A^2-B^2$ exactly when $-AB+BA=0$, or $AB=BA$.
\end{solution}

\begin{problem}\label{ex:01:symmetric-product}
\begin{enumerate}
\item For symmetric $n\times n$ matrices $A,B$, prove that $AB$ is
symmetric exactly when $AB=BA$.
\item Give two symmetric matrices whose product is not symmetric.
\end{enumerate}
\end{problem}
\begin{solution}
(a) Since $A^\top=A$ and $B^\top=B$,
\[
(AB)^\top=B^\top A^\top=BA.
\]
Thus $(AB)^\top=AB$ exactly when $BA=AB$.

(b) The symmetric matrices
$A=\begin{pmatrix}1&0\\0&2\end{pmatrix}$ and
$B=\begin{pmatrix}0&1\\1&0\end{pmatrix}$ give
$AB=\begin{pmatrix}0&1\\2&0\end{pmatrix}$, which is not symmetric.
\end{solution}

\begin{problem}\label{ex:01:symmetric-decomposition}
A real square matrix $K$ is \emph{skew-symmetric} if
$K^\top=-K$. Let $A$ be a real square matrix.
\begin{enumerate}
\item Prove that $A-A^\top$ is skew-symmetric.
\item Express $A$ as the sum of a symmetric matrix and a skew-symmetric
matrix.
\item Prove that every diagonal entry of a skew-symmetric matrix is zero.
\end{enumerate}
\end{problem}
\begin{solution}
(a) Entrywise transposition gives
$(A-A^\top)^\top=A^\top-A=-(A-A^\top)$.

(b) Set
\[
H=\frac{A+A^\top}2,\qquad K=\frac{A-A^\top}2.
\]
Then $A=H+K$, $H^\top=H$, and $K^\top=-K$.

(c) If $K^\top=-K$, the $(i,i)$ entry gives $k_{ii}=-k_{ii}$.
Hence $2k_{ii}=0$, so $k_{ii}=0$.
\end{solution}

\begin{problem}\label{ex:01:outer-product-symmetry}
For real columns $a,b$ of the same length, determine exactly when
$ab^\top=ba^\top$.
\end{problem}
\begin{solution}
If $b=0$, both matrices are zero. Suppose $b\ne0$ and choose an index
$j$ with $b_j\ne0$. Equality of the $(i,j)$ entries gives
$a_i b_j=b_i a_j$ for every $i$, so
\[
a=\frac{a_j}{b_j}b.
\]
Conversely, if $a=kb$ for a real $k$, then
$ab^\top=kbb^\top=ba^\top$.
The exact condition is therefore $b=0$ or $a=kb$ for some real $k$;
this includes $a=0$.
\end{solution}

\begin{problem}\label{ex:01:outer-product-sums}
Let $A=[\,a_1\ \cdots\ a_n\,]$ be an $m\times n$ real matrix.
Prove
\[
AA^\top=\sum_{j=1}^n a_j a_j^\top,
\qquad
(A^\top A)_{ij}=a_i\cdot a_j.
\]
\end{problem}
\begin{solution}
Write the entries of $A$ as $a_{rj}$. The $(r,s)$ entry of the first
product is
\[
(AA^\top)_{rs}=\sum_{j=1}^n a_{rj}a_{sj}
=\sum_{j=1}^n(a_j a_j^\top)_{rs}.
\]
For the second product,
\[
(A^\top A)_{ij}=\sum_{r=1}^m a_{ri}a_{rj}=a_i\cdot a_j.
\]
Both identities follow by equality of entries.
\end{solution}

\begin{problem}\label{ex:01:gram-product-cancellation}
Let $A\in\R^{m\times n}$ and $B,C\in\R^{n\times p}$. Prove:
\begin{enumerate}
\item $A^\top A=0$ if and only if $A=0$;
\item $AB=0$ if and only if $A^\top AB=0$;
\item $AB=AC$ if and only if $A^\top AB=A^\top AC$.
\end{enumerate}
\end{problem}
\begin{solution}
(a) If $A^\top A=0$, then, for every column $j$,
\[
0=(A^\top A)_{jj}=\sum_{i=1}^m a_{ij}^2.
\]
A sum of squares of real numbers is zero only when every summand is
zero. Thus all entries of $A$ vanish. The converse follows by multiplication.

(b) The forward implication follows by multiplying $AB=0$ on the left
by $A^\top$. Conversely, if $A^\top AB=0$, then
\[
(AB)^\top(AB)=B^\top A^\top AB=0.
\]
Part (a), applied to $AB$, gives $AB=0$.

(c) Apply part (b) with $B-C$ in place of $B$:
\[
AB=AC\ \Longleftrightarrow\ A(B-C)=0
\ \Longleftrightarrow\ A^\top A(B-C)=0
\ \Longleftrightarrow\ A^\top AB=A^\top AC.
\]
\end{solution}

\begin{problem}\label{ex:01:trace-product}
For a square matrix $M$, define its \emph{trace} by
$\tr(M)=\sum_i M_{ii}$.
\begin{enumerate}
\item If $A,B$ are real $m\times n$ matrices, prove
\[
\tr(A^\top B)=\tr(BA^\top)
=\tr(AB^\top)=\tr(B^\top A).
\]
\item Prove $\tr(ABC)=\tr(CAB)=\tr(BCA)$ and specify the compatible
sizes of $A,B,C$.
\item Must $\tr(ABC)=\tr(ACB)$ whenever both sides are defined?
\end{enumerate}
\end{problem}
\begin{solution}
(a) If the entries of $A,B$ are $a_{ij},b_{ij}$, then
\[
\tr(A^\top B)=\sum_{j=1}^n\sum_{i=1}^m a_{ij}b_{ij},
\qquad
\tr(BA^\top)=\sum_{i=1}^m\sum_{j=1}^n b_{ij}a_{ij}.
\]
These are equal because finite sums can be reordered and real
multiplication commutes. Interchanging $A$ and $B$ gives the other two
traces, with the same sum.

(b) Take $A$ of size $m\times n$, $B$ of size $n\times p$, and $C$
of size $p\times m$. Then $ABC,CAB,BCA$ are square matrices of sizes
$m,p,n$, respectively. For $U$ of size $r\times s$ and $V$ of size
$s\times r$, part (a) applied to $U^\top,V$ gives
$\tr(UV)=\tr(VU)$. Hence
\[
\tr(ABC)=\tr((AB)C)=\tr(C(AB))=\tr(CAB),
\]
and, applying the same rule to $CA$ and $B$,
$\tr(CAB)=\tr(BCA)$.

(c) No. Take
\[
A=\begin{pmatrix}1&0\\0&0\end{pmatrix},\quad
B=\begin{pmatrix}0&1\\0&0\end{pmatrix},\quad
C=\begin{pmatrix}0&0\\1&0\end{pmatrix}.
\]
Here $BC=A$, whereas $CB=\begin{pmatrix}0&0\\0&1\end{pmatrix}$.
Thus $ABC=A$ and $ACB=0$, giving traces $1$ and $0$.
\end{solution}

\needspace{20\baselineskip}
\subsection*{Inverses and repeated multiplication}

\begin{problem}\label{ex:01:inverse-products-transposes}
Let $A,B$ be invertible real $n\times n$ matrices. Prove that $AB$ and
$A^\top$ are invertible, with
\[
(AB)^{-1}=B^{-1}A^{-1},\qquad
(A^\top)^{-1}=(A^{-1})^\top.
\]
\end{problem}
\begin{solution}
Associativity gives both required products:
\[
(AB)(B^{-1}A^{-1})=A(BB^{-1})A^{-1}=AA^{-1}=I,
\]
\[
(B^{-1}A^{-1})(AB)=B^{-1}(A^{-1}A)B=B^{-1}B=I.
\]
For the transpose, the product rule gives
\[
A^\top(A^{-1})^\top=(A^{-1}A)^\top=I,
\qquad
(A^{-1})^\top A^\top=(AA^{-1})^\top=I.
\]
Each displayed candidate is a two-sided inverse, so uniqueness of the
inverse proves the formulas.
\end{solution}

\begin{problem}\label{ex:01:polynomial-inverses}
Let $A$ be a real $n\times n$ matrix. In each case, prove that the
indicated matrix is invertible by finding its inverse:
\begin{enumerate}
\item $I-A$, if $A^k=0$ for a positive integer $k$;
\item $A$, if $A^2+2A+I=0$;
\item $A$, if $A^2-A+I=0$.
\end{enumerate}
Here $A^0=I$ and $A^j$ means a product of $j$ copies of $A$ for $j\geqslant1$.
\end{problem}
\begin{solution}
(a) Put $S=\sum_{j=0}^{k-1}A^j$. Distributivity and associativity give
\[
(I-A)S=S(I-A)=\sum_{j=0}^{k-1}A^j-\sum_{j=1}^{k}A^j
=I-A^k=I.
\]
Thus $(I-A)^{-1}=S$. When $k=1$, the hypothesis says $A=0$ and
$S=I$, so the calculation includes that case.

(b) The hypothesis gives
\[
A(-A-2I)=(-A-2I)A=-A^2-2A=I.
\]
Hence $A^{-1}=-A-2I$.

(c) The hypothesis gives
\[
A(I-A)=(I-A)A=A-A^2=I,
\]
so $A^{-1}=I-A$.
\end{solution}

\begin{problem}\label{ex:01:centering-matrix}
Let $n\geqslant1$, let $j=(1,\ldots,1)^\top\in\R^n$, and set
\[
P=I_n-\frac1n jj^\top.
\]
Prove that $P$ is symmetric and \emph{idempotent}, meaning $P^2=P$.
Explain these properties by describing $Px$ for a column $x\in\R^n$.
\end{problem}
\begin{solution}
Since $j^\top j=n$, associativity gives
\[
(jj^\top)^2=j(j^\top j)j^\top=n jj^\top.
\]
Therefore
\[
P^2=I_n-\frac2n jj^\top+\frac1{n^2}(jj^\top)^2
=I_n-\frac1n jj^\top=P.
\]
Also $(jj^\top)^\top=jj^\top$, so
$P^\top=P$.
If $\bar x=\frac1n\sum_{i=1}^n x_i$, then
\[
Px=x-\bar xj.
\]
Thus $P$ subtracts the mean from each coordinate. The result has
coordinate sum $\sum_i(x_i-\bar x)=0$, so applying $P$ again changes
nothing. Its diagonal entries are $1-1/n$ and its off-diagonal entries
are $-1/n$: each coordinate contributes the same amount to every other
coordinate, giving symmetry.
\end{solution}

\needspace{20\baselineskip}
\subsection*{Homogeneous systems and solution sets}

\begin{problem}\label{ex:01:two-equations-three-unknowns}
For each system, write its coefficient matrix, reduce it to row echelon
form, and find all solutions.
\begin{enumerate}[label=(\alph*)]
\item $2x-3y+4z=0$, $3x+y+z=0$.
\item $-x+2y-4z+w=0$, $x+3y+z-w=0$.
\end{enumerate}
\end{problem}
\begin{solution}
(a) The coefficient matrix and a row reduction are
\[
\begin{pmatrix}2&-3&4\\3&1&1\end{pmatrix}
\xrightarrow{R_2\leftarrow R_2-\frac32R_1}
\begin{pmatrix}2&-3&4\\0&11/2&-5\end{pmatrix}.
\]
The second row gives $y=10z/11$. Substituting into the first gives
\[
2x=3y-4z=\frac{30}{11}z-4z=-\frac{14}{11}z,
\qquad x=-\frac7{11}z.
\]
Writing $z=11t$, all solutions are $t(-7,10,11)^\top$, $t\in\R$.

(b) Here
\[
\begin{pmatrix}-1&2&-4&1\\1&3&1&-1\end{pmatrix}
\xrightarrow{R_2\leftarrow R_2+R_1}
\begin{pmatrix}-1&2&-4&1\\0&5&-3&0\end{pmatrix}.
\]
Thus $y=3z/5$, and the first row gives
$x=2y-4z+w=6z/5-4z+w=w-14z/5$.
With $z=5s$ and $w=t$, all solutions are
\[
(x,y,z,w)^\top=(-14s+t,3s,5s,t)^\top,\qquad s,t\in\R.
\]
In both parts the free coordinates can be chosen arbitrarily, and back
substitution determines the remaining coordinates. Hence the displayed
sets include every solution; substituting the formulas verifies them.
\end{solution}

\begin{problem}\label{ex:01:combining-homogeneous-solutions}
Let $a_1,\ldots,a_m\in\R^n$. Suppose $x\cdot a_i=y\cdot a_i=0$
for every $i$. Prove that $x+y$ also satisfies these equations and that
$cx$ satisfies them for every real number $c$.
\end{problem}
\begin{solution}
For each $i$, distributivity gives
$(x+y)\cdot a_i=x\cdot a_i+y\cdot a_i=0$.
Scalar compatibility gives $(cx)\cdot a_i=c(x\cdot a_i)=0$.
These equalities hold for every row, proving both assertions.
\end{solution}

\begin{problem}\label{ex:01:translated-solutions}
Let $A\in\R^{m\times n}$ and suppose $Ax=b$.
If $Ax'=b$ as well, prove that $x'=x+y$ for some $y$ satisfying $Ay=0$.
Conversely, show that $x+y$ solves the original system whenever $Ay=0$.
\end{problem}
\begin{solution}
Set $y=x'-x$. Then $Ay=Ax'-Ax=b-b=0$ and $x'=x+y$.
Conversely, $A(x+y)=Ax+Ay=b+0=b$. Thus the full solution set is
one particular solution plus all solutions of the homogeneous system.
\end{solution}

\begin{problem}\label{ex:01:affine-combinations-of-solutions}
Let $A\in\R^{m\times n}$, $B\in\R^{m\times p}$, and suppose
$AX_i=B$ for $i=1,\ldots,k$, where every $X_i$ is $n\times p$.
If real numbers $c_1,\ldots,c_k$ satisfy $\sum_i c_i=1$, prove that
$\sum_i c_iX_i$ also solves $AX=B$.
\end{problem}
\begin{solution}
All matrices in the sum have the same size. Distributivity gives
\[
A\left(\sum_{i=1}^k c_iX_i\right)
=\sum_{i=1}^k c_iAX_i
=\sum_{i=1}^k c_iB
=\left(\sum_{i=1}^k c_i\right)B=B.
\]
Thus the asserted matrix is a solution, including when $B=0$.
\end{solution}

\needspace{20\baselineskip}
\subsection*{Consistency and exceptional parameters}

\begin{problem}\label{ex:01:dependent-equations-parameter}
For which real numbers $k$ does the system
\[
x_1+x_2=1,\qquad 2x_1+2x_2=k
\]
have a solution? When it has a solution, is that solution unique?
\end{problem}
\begin{solution}
Subtracting twice the first equation from the second gives $0=k-2$.
Thus a solution exists exactly when $k=2$. In that case the second
equation repeats the first, and all solutions are $(1-t,t)^\top$, $t\in\R$.
There are infinitely many, so the solution is not unique.
\end{solution}

\begin{problem}\label{ex:01:universal-right-hand-sides}
Can a matrix $A$ be chosen so that $Ax=b$ is inconsistent for every
right-hand side $b$ of the appropriate size?
Can a matrix $A$ be chosen so that $Ax=b$ is consistent for every such
$b$? Justify each answer.
\end{problem}
\begin{solution}
The first answer is no: for every matrix $A$, the right-hand side $b=0$
has the solution $x=0$.
The second answer is yes: choose $A=I_m$. Then $x=b$ solves $Ax=b$
for every $b\in\R^m$.
\end{solution}

\begin{problem}\label{ex:01:inconsistent-three-equations}
Does the following system have a solution?
\[
\begin{aligned}
2x_1+3x_2+x_3&=7,\\
x_1+2x_2-2x_3&=8,\\
3x_1+4x_2+4x_3&=20.
\end{aligned}
\]
\end{problem}
\begin{solution}
The coefficient rows satisfy
\[
2(2,3,1)-(1,2,-2)=(3,4,4).
\]
Thus twice the first left-hand side minus the second equals the third
left-hand side. The same operation on the right-hand sides gives
$2(7)-8=6$, not $20$. Subtracting twice row one and adding row two to
row three therefore produces $0=14$. The system is inconsistent.
\end{solution}

\begin{problem}\label{ex:01:two-parameter-consistency}
For real parameters $\alpha,\beta$, consider
\[
\begin{aligned}
\alpha x_1+\beta x_2+2x_3&=1,\\
\alpha x_1+(2\beta-1)x_2+3x_3&=1,\\
\alpha x_1+\beta x_2+(\beta+3)x_3&=2\beta-1.
\end{aligned}
\]
When does a solution exist, and when is it unique?
\end{problem}
\begin{solution}
Subtract the first equation from the second and third. The equivalent
system is
\[
\alpha x_1+\beta x_2+2x_3=1,\qquad
(\beta-1)x_2+x_3=0,\qquad
(\beta+1)x_3=2(\beta-1).
\]
If $\beta=-1$, the last equation is $0=-4$, so there is no solution.
If $\beta=1$, then $x_3=0$ and $\alpha x_1+x_2=1$; all solutions are
$(t,1-\alpha t,0)^\top$, $t\in\R$, for every $\alpha$.

For $\beta\ne\pm1$, the last equation first gives $x_3$, and the
second gives $x_2=-x_3/(\beta-1)$. Hence
\[
x_3=\frac{2(\beta-1)}{\beta+1},\qquad
x_2=-\frac2{\beta+1}.
\]
Substituting into the first equation now gives
\[
\alpha x_1=1-\beta x_2-2x_3
=1+\frac{2\beta}{\beta+1}-\frac{4(\beta-1)}{\beta+1}
=\frac{5-\beta}{\beta+1}.
\]
If $\alpha\ne0$, there is exactly one solution, with
$x_1=(5-\beta)/(\alpha(\beta+1))$.
If $\alpha=0$, consistency requires $\beta=5$; the solutions are then
$(t,-1/3,4/3)^\top$, $t\in\R$.
Thus uniqueness holds exactly when $\alpha\ne0$ and $\beta\ne\pm1$.
The remaining consistent cases are $\beta=1$ with arbitrary $\alpha$,
and $(\alpha,\beta)=(0,5)$; each has infinitely many solutions.
All other parameter pairs give no solution.
\end{solution}

\needspace{20\baselineskip}
\subsection*{Point--plane distances}

\begin{problem}\label{ex:01:point-plane-distance}
Find the distance between each point and plane:
\begin{enumerate}
\item $(1,1,2)$ and $3x+y-5z=2$;
\item $(-1,3,2)$ and $2x-4y+z=1$;
\item $(3,-2,1)$ and the $yz$-plane;
\item $(-3,-2,1)$ and the $yz$-plane.
\end{enumerate}
\end{problem}
\begin{solution}
For a plane $n\cdot x=c$ with $n\ne0$ and a point $q$, the distance is
$|n\cdot q-c|/\norm n$.
(a) Here $n=(3,1,-5)$, $c=2$, and $\norm n=\sqrt{9+1+25}=\sqrt{35}$.
Thus
\[
d=\frac{|3(1)+1-5(2)-2|}{\sqrt{35}}=\frac8{\sqrt{35}}.
\]
(b) Here $n=(2,-4,1)$, $c=1$, and $\norm n=\sqrt{4+16+1}=\sqrt{21}$.
Hence
\[
d=\frac{|2(-1)-4(3)+2-1|}{\sqrt{21}}=\frac{13}{\sqrt{21}}.
\]
(c,d) The $yz$-plane has equation $x=0$, with unit normal $(1,0,0)$.
The two distances are therefore $|3|=3$ and $|-3|=3$.
The points lie on opposite sides of the plane; the absolute value makes
both distances nonnegative.
\end{solution}
